\chapter{Fazit}
\label{ch:fazit}

Die Arbeit untersuchte, wie sich verschiedene Zephyr-Konfigurationen auf das Laufzeitverhalten und den Speicherbedarf eines ressourcenbeschr{\"a}nkten Mikrocontrollers auswirken. Daf{\"u}r wurden ein Minimalprofil, ein Logging-Profil und ein IoT-Profil mit OpenThread, CoAP und DTLS/ PSK auf demselben nRF52840 DK verglichen.

Die Ergebnisse zeigen, dass weitere Subsysteme vor allem den statischen Speicherbedarf beeinflussen. Logging braucht im Vergleich zu Minimal 9304 Byte mehr Flash und 2120 Byte mehr RAM. Beim IoT-Profil steigen diese Werte um 165.612 Byte Flash und 71.448 Byte RAM. Die gemessenen Kontextwechselzeiten erh{\"o}hen sich dabei nur um ungef{\"a}hr 5 bis 8 \%. Heap-Operationen ver{\"a}ndern sich nur gering, w{\"a}hrend die Interrupt-Ergebnisse keine monotone Entwicklung zeigen. Beim Timer-Jitter war innerhalb der verf{\"u}gbaren RTC-Aufl{\"o}sung kein Profilunterschied sichtbar.

Die Forschungsfrage l{\"a}sst sich damit folgenderma{\ss}en beantworten: Der Einfluss der Zephyr-Konfiguration zeigt sich auf der untersuchten Hardware besonders deutlich beim Speicherbedarf. Der Einfluss auf kurze, isolierte Kerneloperationen ist hierbei geringer und pfadabh{\"a}ngig und kann von {\"A}nderungen des Bin{\"a}rlayouts {\"u}berlagert werden. Eine Konfiguration sollte daher nicht anhand einer einzelnen Laufzeitmetrik beurteilt werden. F{\"u}r ressourcenbeschr{\"a}nkte Systeme m{\"u}ssen Funktionsumfang, Heap- und Stackreservierungen, reale Laufzeitlast und die verf{\"u}gbare Messaufl{\"o}sung zusammen betrachtet werden.
